\documentclass[12pt]{article}
\usepackage[T1]{fontenc}
\usepackage[utf8]{inputenc}
\usepackage{lmodern}
\usepackage{textcomp}
\usepackage{enumitem}
\usepackage{geometry}
\usepackage{graphicx}
\usepackage{amsmath, amssymb}
\geometry{margin=1in}
\begin{document}
\begin{center}
Psychology - Graduate Level Assessment 12 \
Total Marks: 40 \
Answer all questions.
\end{center}

\section*{Multiple Choice (10 marks)}
\begin{enumerate}[label=\arabic*.]
\item We know less about why certain stimuli are attractive to humansthan we do about what is attractive to fish. Discuss thisstatement.
\begin{enumerate}[label=(\alph*)]
    \item Attractiveness can be precisely measured and quantified in all individuals.
    \item We fully understand the factors that determine human attractiveness.
    \item Attractiveness is a trait that humans are born understanding completely.
    \item Human attractiveness is determined by a single genetic marker.
    \item The principles of attractiveness in humans have been fully mapped out by neuroscience.
\end{enumerate}
\item Piaget's theory describes stages of cognitive development \_\_\_\_\_\_\_\_\_\_\_.
\begin{enumerate}[label=(\alph*)]
    \item During adolescence only
    \item From birth to late childhood only
    \item During infancy only
    \item From birth to early adulthood only
    \item From infancy through adolescence only
\end{enumerate}
\item The belief that a child's misbehavior has one of four goals - i.e., attention, revenge, power, or to display inadequacy - is most consistent with:
\begin{enumerate}[label=(\alph*)]
    \item Bandura's social learning theory.
    \item Skinner's operant conditioning theory.
    \item Maslow's hierarchy of needs theory.
    \item Beck's cognitive-behavioral therapy.
    \item Pavlov's classical conditioning theory.
\end{enumerate}
\item What do psychologists generally accept as the causes of functional psychoses?
\begin{enumerate}[label=(\alph*)]
    \item Childhood experiences
    \item Astrological signs and birth months
    \item Educational background and intelligence quotient (IQ)
    \item Exercise frequency and body mass index (BMI)
    \item Use of technology and media consumption
\end{enumerate}
\item \_\_\_\_\_\_\_\_\_\_ is the most frequent cause of many of the ethical complaints brought against psychotherapy supervisors.
\begin{enumerate}[label=(\alph*)]
    \item Negligence in supervision
    \item Incompetence
    \item Misuse of authority
    \item Sexual harassment
    \item Lack of timely feedback
\end{enumerate}
\end{enumerate}

\section*{True / False (10 marks)}
\textit{Answer True or False}
\begin{enumerate}[label=\arabic*.]
\setcounter{enumi}{5}
\item True or False: Surprise is not classified as one of the primary facial expressions recognized in psychological research.
\item True or False: Time-out procedures rely on classical conditioning to associate negative experiences with specific behaviors.
\item True or False: Stimulus satiation involves the repeated presentation of a stimulus until its attractiveness is diminished for the client.
\item True or False: Strong organizational culture is correlated with low levels of job commitment in the workplace.
\item True or False: The Whorfian hypothesis asserts that all languages share the same phonetic structures, leading to universal tone systems.
\end{enumerate}

\section*{Long Form Response (20 marks)}
\textit{Answer each question with a clear and complete explanation.}
\begin{enumerate}[label=\arabic*.]
\setcounter{enumi}{10}
\item Discuss the concept of Type II error and the factors that influence statistical power, specifically analyzing the impact of sample size on power in hypothesis testing.
\item Discuss the role of the basal ganglia in the Cannon-Bard theory of emotion, analyzing its significance and implications for understanding emotional processing in the brain.
\end{enumerate}

\vfill
\noindent\textit{End of Paper}
\end{document}